\chapter{Fourier Transform and Frequency Analysis}
\label{ch:fourier}
% ============================================================================

\section{Fourier Series}

Any periodic function $f(t)$ with period $T$ can be expressed as:
\begin{equation}
f(t) = a_0 + \sum_{n=1}^{\infty} \left[ a_n \cos\left(\frac{2\pi nt}{T}\right) + b_n \sin\left(\frac{2\pi nt}{T}\right) \right]
\end{equation}

where:
\begin{align}
a_0 &= \frac{1}{T} \int_0^T f(t) \, \dd t \\
a_n &= \frac{2}{T} \int_0^T f(t) \cos\left(\frac{2\pi nt}{T}\right) \dd t \\
b_n &= \frac{2}{T} \int_0^T f(t) \sin\left(\frac{2\pi nt}{T}\right) \dd t
\end{align}

Complex exponential form:
\begin{equation}
f(t) = \sum_{n=-\infty}^{\infty} c_n \ee^{\jj n \omega_0 t}, \quad \omega_0 = \frac{2\pi}{T}
\end{equation}

where $c_n = \frac{1}{T} \int_0^T f(t) \ee^{-\jj n \omega_0 t} \, \dd t$

\begin{figure}[htbp]
\centering
\textbf{Fourier Series Decomposition of a Square Wave}

\vspace{0.5cm}

\begin{tikzpicture}
\begin{axis}[
    width=13cm, height=5cm,
    xlabel={Time $t$},
    ylabel={Amplitude},
    xmin=-1.1, xmax=1.1,
    ymin=-1.5, ymax=1.5,
    xtick={-1, -0.5, 0, 0.5, 1},
    xticklabels={$-T$, $-T/2$, $0$, $T/2$, $T$},
    ytick={-1, 0, 1},
    grid=major,
    grid style={gray!30},
    legend pos=outer north east,
    legend style={font=\footnotesize},
    title={\normalsize Individual Harmonic Components},
    title style={at={(0.5,1.05)}, anchor=south}
]

% DC component (a_0 = 0 for symmetric square wave)

% Fundamental (n=1): (4/pi)*sin(pi*x)
\addplot[UofSCGarnet, very thick, domain=-1.1:1.1, samples=200]
    {(4/pi)*sin(deg(pi*x))};
\addlegendentry{$n=1$: $(4/\pi)\sin(\omega_0 t)$}

% 3rd harmonic: (4/3pi)*sin(3*pi*x)
\addplot[UofSCAtlantic, thick, domain=-1.1:1.1, samples=200]
    {(4/(3*pi))*sin(deg(3*pi*x))};
\addlegendentry{$n=3$: $(4/3\pi)\sin(3\omega_0 t)$}

% 5th harmonic: (4/5pi)*sin(5*pi*x)
\addplot[UofSCHorseshoe, thick, domain=-1.1:1.1, samples=200]
    {(4/(5*pi))*sin(deg(5*pi*x))};
\addlegendentry{$n=5$: $(4/5\pi)\sin(5\omega_0 t)$}

% 7th harmonic
\addplot[UofSCCongaree, thick, domain=-1.1:1.1, samples=200]
    {(4/(7*pi))*sin(deg(7*pi*x))};
\addlegendentry{$n=7$: $(4/7\pi)\sin(7\omega_0 t)$}

\end{axis}
\end{tikzpicture}

\vspace{0.8cm}

\begin{tikzpicture}
\begin{axis}[
    width=13cm, height=5cm,
    xlabel={Time $t$},
    ylabel={Amplitude},
    xmin=-1.1, xmax=1.1,
    ymin=-1.5, ymax=1.5,
    xtick={-1, -0.5, 0, 0.5, 1},
    xticklabels={$-T$, $-T/2$, $0$, $T/2$, $T$},
    ytick={-1, 0, 1},
    grid=major,
    grid style={gray!30},
    legend pos=outer north east,
    legend style={font=\footnotesize},
    title={\normalsize Cumulative Sum of Harmonics},
    title style={at={(0.5,1.05)}, anchor=south}
]

% Ideal square wave (for reference)
\addplot[UofSC50Black, dashed, very thick, domain=-1.1:-0.5, samples=2, forget plot] {-1};
\addplot[UofSC50Black, dashed, very thick, domain=-0.5:-0.5, samples=2, forget plot] coordinates {(-0.5,-1) (-0.5,1)};
\addplot[UofSC50Black, dashed, very thick, domain=-0.5:0.5, samples=2, forget plot] {1};
\addplot[UofSC50Black, dashed, very thick, domain=0.5:0.5, samples=2, forget plot] coordinates {(0.5,1) (0.5,-1)};
\addplot[UofSC50Black, dashed, very thick, domain=0.5:1.1, samples=2] {-1};
\addlegendentry{Ideal Square Wave}

% n=1 only
\addplot[UofSCGarnet, thick, domain=-1.1:1.1, samples=200]
    {(4/pi)*sin(deg(pi*x))};
\addlegendentry{$N=1$}

% n=1,3
\addplot[UofSCAtlantic, thick, domain=-1.1:1.1, samples=200]
    {(4/pi)*sin(deg(pi*x)) + (4/(3*pi))*sin(deg(3*pi*x))};
\addlegendentry{$N=3$}

% n=1,3,5
\addplot[UofSCHorseshoe, thick, domain=-1.1:1.1, samples=200]
    {(4/pi)*sin(deg(pi*x)) + (4/(3*pi))*sin(deg(3*pi*x)) + (4/(5*pi))*sin(deg(5*pi*x))};
\addlegendentry{$N=5$}

% n=1,3,5,7,9,11
\addplot[UofSCCongaree, very thick, domain=-1.1:1.1, samples=300]
    {(4/pi)*sin(deg(pi*x)) + (4/(3*pi))*sin(deg(3*pi*x)) + (4/(5*pi))*sin(deg(5*pi*x)) + (4/(7*pi))*sin(deg(7*pi*x)) + (4/(9*pi))*sin(deg(9*pi*x)) + (4/(11*pi))*sin(deg(11*pi*x))};
\addlegendentry{$N=11$}

% n=1,3,5,7,9,11,13,15,17,19
\addplot[UofSCRose, very thick, domain=-1.1:1.1, samples=400]
    {(4/pi)*sin(deg(pi*x)) + (4/(3*pi))*sin(deg(3*pi*x)) + (4/(5*pi))*sin(deg(5*pi*x)) + (4/(7*pi))*sin(deg(7*pi*x)) + (4/(9*pi))*sin(deg(9*pi*x)) + (4/(11*pi))*sin(deg(11*pi*x)) + (4/(13*pi))*sin(deg(13*pi*x)) + (4/(15*pi))*sin(deg(15*pi*x)) + (4/(17*pi))*sin(deg(17*pi*x)) + (4/(19*pi))*sin(deg(19*pi*x))};
\addlegendentry{$N=19$}

\end{axis}
\end{tikzpicture}

\caption{Fourier series decomposition of a square wave. \textbf{Top:} Individual odd harmonic components ($n = 1, 3, 5, 7$) with amplitudes $4/(n\pi)$. A symmetric square wave contains only odd harmonics. \textbf{Bottom:} Cumulative sum showing how adding more harmonics improves the approximation. The dashed gray line shows the ideal square wave. Note the Gibbs phenomenon (overshoot near discontinuities) which persists regardless of how many terms are included. The Fourier series formula is: $f(t) = \sum_{n=1,3,5,\ldots}^{\infty} \frac{4}{n\pi}\sin(n\omega_0 t)$ where $\omega_0 = 2\pi/T$.}
\label{fig:fourier_series}
\end{figure}


\begin{examplebox}[title=Example: Fourier Series of a Square Wave]
Consider a square wave with period $T = 2\pi$ and amplitude $A$:
\[
f(t) = \begin{cases}
A & 0 < t < \pi \\
-A & \pi < t < 2\pi
\end{cases}
\]

\textbf{Step 1: Compute $a_0$}
\[
a_0 = \frac{1}{2\pi} \int_0^{2\pi} f(t) \, \dd t = \frac{1}{2\pi} \left[ \int_0^{\pi} A \, \dd t + \int_{\pi}^{2\pi} (-A) \, \dd t \right] = \frac{1}{2\pi}(A\pi - A\pi) = 0
\]

\textbf{Step 2: Compute $a_n$}

Since $f(t)$ is an odd function about $t = \pi$, all cosine coefficients vanish: $a_n = 0$ for all $n$.

\textbf{Step 3: Compute $b_n$}
\begin{align*}
b_n &= \frac{2}{2\pi} \int_0^{2\pi} f(t) \sin(nt) \, \dd t \\
&= \frac{1}{\pi} \left[ \int_0^{\pi} A\sin(nt) \, \dd t + \int_{\pi}^{2\pi} (-A)\sin(nt) \, \dd t \right] \\
&= \frac{A}{\pi} \left[ -\frac{\cos(nt)}{n} \Big|_0^{\pi} + \frac{\cos(nt)}{n} \Big|_{\pi}^{2\pi} \right] \\
&= \frac{A}{n\pi} \left[ (-\cos(n\pi) + 1) + (\cos(2n\pi) - \cos(n\pi)) \right] \\
&= \frac{A}{n\pi} \left[ 2 - 2\cos(n\pi) \right] = \frac{2A}{n\pi}(1 - (-1)^n)
\end{align*}

For $n$ odd: $b_n = \dfrac{4A}{n\pi}$. For $n$ even: $b_n = 0$.

\textbf{Result:} The Fourier series of the square wave is:
\begin{equation}
\boxed{f(t) = \frac{4A}{\pi} \left( \sin(t) + \frac{1}{3}\sin(3t) + \frac{1}{5}\sin(5t) + \cdots \right) = \frac{4A}{\pi} \sum_{n=1,3,5,\ldots}^{\infty} \frac{\sin(nt)}{n}}
\end{equation}

This shows that a square wave contains only odd harmonics, with amplitudes decreasing as $1/n$.
\end{examplebox}

\begin{keyconceptbox}[title=Gibbs Phenomenon]
When approximating a discontinuous function with a finite Fourier series, overshoot occurs near discontinuities. This overshoot is approximately 9\% of the jump magnitude and does not diminish as more terms are added---it only becomes narrower. This is known as the \textbf{Gibbs phenomenon}.
\end{keyconceptbox}

\section{Fourier Transform}

For non-periodic functions, the Fourier transform extends the series to continuous frequencies:

\begin{definitionbox}[title=Definition 5.1: Fourier Transform]
\begin{align}
F(\omega) &= \Fourier\{f(t)\} = \int_{-\infty}^{\infty} f(t) \ee^{-\jj\omega t} \, \dd t \\
f(t) &= \iFourier\{F(\omega)\} = \frac{1}{2\pi} \int_{-\infty}^{\infty} F(\omega) \ee^{\jj\omega t} \, \dd\omega
\end{align}
\end{definitionbox}

\begin{figure}[htbp]
\centering
\textbf{Frequency Spectrum of a Damped Sinusoid}

\vspace{0.3cm}

% Time domain signal
\begin{tikzpicture}
\begin{axis}[
    width=13cm, height=4.5cm,
    xlabel={Time $t$ (seconds)},
    ylabel={$f(t)$},
    xmin=0, xmax=5,
    ymin=-1.1, ymax=1.1,
    grid=major,
    grid style={gray!30},
    title={\normalsize Time Domain: $f(t) = \ee^{-\sigma t}\cos(\omega_0 t)$ \quad ($\sigma = 0.5$, $\omega_0 = 4\pi$ rad/s)},
    title style={at={(0.5,1.05)}, anchor=south}
]

% Damped sinusoid: exp(-0.5*t)*cos(4*pi*t) = exp(-0.5*t)*cos(2*pi*2*t), f_0 = 2 Hz
\addplot[UofSCAtlantic, very thick, domain=0:5, samples=500]
    {exp(-0.5*x)*cos(deg(4*pi*x))};

% Envelope (positive)
\addplot[UofSCGarnet, dashed, thick, domain=0:5, samples=100]
    {exp(-0.5*x)};

% Envelope (negative)
\addplot[UofSCGarnet, dashed, thick, domain=0:5, samples=100]
    {-exp(-0.5*x)};

% Annotation
\node[UofSCGarnet, font=\small] at (3.2, 0.35) {Envelope: $\ee^{-\sigma t}$};
\draw[UofSCGarnet, thin, ->] (2.7, 0.3) -- (2.0, 0.38);

\node[UofSCAtlantic, font=\small] at (1.5, -0.65) {$f_0 = 2$ Hz};

\end{axis}
\end{tikzpicture}

\vspace{0.6cm}

% Magnitude spectrum
\begin{tikzpicture}
\begin{axis}[
    width=6.2cm, height=5cm,
    xlabel={Frequency $f$ (Hz)},
    ylabel={$|F(\jj\omega)|$},
    xmin=-6, xmax=6,
    ymin=0, ymax=1.2,
    grid=major,
    grid style={gray!30},
    title={\normalsize Magnitude Spectrum},
    title style={at={(0.5,1.05)}, anchor=south},
    xtick={-6,-4,-2,0,2,4,6},
    ytick={0, 0.5, 1}
]

% Magnitude spectrum: |F(jw)| = 0.5 * [sigma/((sigma)^2 + (w-w0)^2) + sigma/((sigma)^2 + (w+w0)^2)]
% For damped cosine: peaks at +/- f_0 with width determined by sigma
% Using Lorentzian-like shape centered at f0=2 Hz

% Peak at +f_0 = 2 Hz
\addplot[UofSCAtlantic, very thick, domain=-6:6, samples=300]
    {0.5*0.5/((0.5)^2 + (2*pi*(x-2))^2) + 0.5*0.5/((0.5)^2 + (2*pi*(x+2))^2)};

% Mark the center frequencies
\draw[UofSCGarnet, dashed, thick] (2, 0) -- (2, 1.1);
\draw[UofSCGarnet, dashed, thick] (-2, 0) -- (-2, 1.1);

% Annotations
\node[UofSCGarnet, font=\footnotesize, anchor=south] at (2, 1.12) {$f_0$};
\node[UofSCGarnet, font=\footnotesize, anchor=south] at (-2, 1.12) {$-f_0$};

% Half-power bandwidth annotation
\draw[UofSCHorseshoe, thick, <->] (1.84, 0.5) -- (2.16, 0.5);
\node[UofSCHorseshoe, font=\tiny, anchor=west] at (2.2, 0.5) {$\Delta f$};

\end{axis}
\end{tikzpicture}
\hfill
% Phase spectrum
\begin{tikzpicture}
\begin{axis}[
    width=6.2cm, height=5cm,
    xlabel={Frequency $f$ (Hz)},
    ylabel={$\angle F(\jj\omega)$ (degrees)},
    xmin=-6, xmax=6,
    ymin=-100, ymax=100,
    grid=major,
    grid style={gray!30},
    title={\normalsize Phase Spectrum},
    title style={at={(0.5,1.05)}, anchor=south},
    xtick={-6,-4,-2,0,2,4,6},
    ytick={-90, -45, 0, 45, 90},
    yticklabels={$-90$, $-45$, $0$, $45$, $90$}
]

% Phase spectrum: transitions through the resonance frequencies
% Phase = atan2(Im, Re) for the transform
% Near +f_0: phase transitions from +90 to -90
% Near -f_0: phase transitions from -90 to +90

% Approximate phase behavior (simplified for visualization)
\addplot[UofSCCongaree, very thick, domain=-6:-2.3, samples=50] {-90 + 90*2/pi*atan(deg(5*(x+2)))};
\addplot[UofSCCongaree, very thick, domain=-1.7:1.7, samples=50] {0};
\addplot[UofSCCongaree, very thick, domain=1.7:6, samples=50] {90 - 90*2/pi*atan(deg(5*(x-2)))};

% Transition regions
\addplot[UofSCCongaree, very thick, domain=-2.3:-1.7, samples=30]
    {-90*tanh(3*(x+2))};
\addplot[UofSCCongaree, very thick, domain=1.7:2.3, samples=30]
    {-90*tanh(3*(x-2))};

% Mark center frequencies
\draw[UofSCGarnet, dashed, thick] (2, -100) -- (2, 100);
\draw[UofSCGarnet, dashed, thick] (-2, -100) -- (-2, 100);

% Annotations
\node[UofSCGarnet, font=\footnotesize, anchor=south] at (2, 92) {$f_0$};
\node[UofSCGarnet, font=\footnotesize, anchor=south] at (-2, 92) {$-f_0$};

\end{axis}
\end{tikzpicture}

\vspace{0.4cm}

% Discrete spectrum example (for periodic signal)
\begin{tikzpicture}
\begin{axis}[
    width=13cm, height=4cm,
    xlabel={Harmonic Number $n$},
    ylabel={$|c_n|$},
    xmin=-0.5, xmax=12.5,
    ymin=0, ymax=1.4,
    grid=major,
    grid style={gray!30},
    title={\normalsize Discrete Magnitude Spectrum of Square Wave (Fourier Series Coefficients)},
    title style={at={(0.5,1.05)}, anchor=south},
    xtick={0,1,2,3,4,5,6,7,8,9,10,11,12},
    ytick={0, 0.5, 1, 1.27},
    yticklabels={$0$, $0.5$, $1$, $4/\pi$}
]

% Magnitude of Fourier coefficients for square wave: |c_n| = 4/(n*pi) for odd n, 0 for even n
% Plot as stem plot using ycomb

\addplot+[UofSCAtlantic, very thick, ycomb, mark=*, mark options={fill=UofSCAtlantic, scale=1.2}]
    coordinates {
        (0, 0)
        (1, 1.273)
        (2, 0)
        (3, 0.424)
        (4, 0)
        (5, 0.255)
        (6, 0)
        (7, 0.182)
        (8, 0)
        (9, 0.141)
        (10, 0)
        (11, 0.116)
        (12, 0)
    };

% Envelope curve (1/n decay)
\addplot[UofSCGarnet, dashed, thick, domain=0.8:12, samples=50]
    {4/(pi*x)};

% Annotations
\node[UofSCGarnet, font=\footnotesize] at (8, 0.65) {Envelope: $4/(n\pi)$};
\node[UofSCAtlantic, font=\footnotesize] at (6.5, 0.12) {Even harmonics = 0};

\end{axis}
\end{tikzpicture}

\caption{Frequency domain representation of signals. \textbf{Top:} Time-domain damped sinusoid $f(t) = \ee^{-\sigma t}\cos(\omega_0 t)$ with decay rate $\sigma = 0.5$ s$^{-1}$ and frequency $f_0 = 2$ Hz. The dashed Garnet lines show the exponential envelope. \textbf{Middle:} Magnitude spectrum (Atlantic blue) shows peaks at $\pm f_0$ with bandwidth $\Delta f$ proportional to $\sigma$. Phase spectrum (Congaree teal) shows $\pm 90°$ phase transitions at the resonance frequencies. \textbf{Bottom:} Discrete spectrum for a periodic square wave showing only odd harmonics with magnitude decreasing as $4/(n\pi)$.}
\label{fig:spectrum_visualization}
\end{figure}


\begin{examplebox}[title=Example: Fourier Transform of a Decaying Exponential]
Find the Fourier transform of $f(t) = \ee^{-at}u(t)$ where $a > 0$ and $u(t)$ is the unit step function.

\textbf{Solution:}
\begin{align*}
F(\omega) &= \int_{-\infty}^{\infty} \ee^{-at}u(t) \ee^{-\jj\omega t} \, \dd t \\
&= \int_{0}^{\infty} \ee^{-at} \ee^{-\jj\omega t} \, \dd t \\
&= \int_{0}^{\infty} \ee^{-(a + \jj\omega)t} \, \dd t \\
&= \left[ \frac{\ee^{-(a + \jj\omega)t}}{-(a + \jj\omega)} \right]_{0}^{\infty}
\end{align*}

Since $a > 0$, the exponential vanishes as $t \to \infty$:
\[
F(\omega) = 0 - \frac{1}{-(a + \jj\omega)} = \frac{1}{a + \jj\omega}
\]

\textbf{Result:}
\begin{equation}
\boxed{\Fourier\{\ee^{-at}u(t)\} = \frac{1}{a + \jj\omega}}
\end{equation}

\textbf{Magnitude and Phase:}
\[
|F(\omega)| = \frac{1}{\sqrt{a^2 + \omega^2}}, \quad \angle F(\omega) = -\arctan\left(\frac{\omega}{a}\right)
\]

This result shows that a decaying exponential has a low-pass frequency spectrum: low frequencies pass through with little attenuation, while high frequencies are attenuated.
\end{examplebox}

\subsection{Fourier Transform Properties}

The Fourier transform possesses several important properties that make it a powerful tool for signal analysis:

\begin{table}[htbp]
\centering
\caption{Fourier Transform Properties}
\label{tab:fourier_properties}
\begin{tabular}{@{}lcc@{}}
\toprule
\textbf{Property} & \textbf{Time Domain} & \textbf{Frequency Domain} \\
\midrule
Linearity & $af(t) + bg(t)$ & $aF(\omega) + bG(\omega)$ \\[6pt]
Time Shifting & $f(t - t_0)$ & $F(\omega)\ee^{-\jj\omega t_0}$ \\[6pt]
Frequency Shifting & $f(t)\ee^{\jj\omega_0 t}$ & $F(\omega - \omega_0)$ \\[6pt]
Time Scaling & $f(at)$ & $\dfrac{1}{|a|}F\left(\dfrac{\omega}{a}\right)$ \\[6pt]
Time Reversal & $f(-t)$ & $F(-\omega)$ \\[6pt]
Differentiation & $\dfrac{\dd f}{\dd t}$ & $\jj\omega F(\omega)$ \\[6pt]
Integration & $\displaystyle\int_{-\infty}^{t} f(\tau)\,\dd\tau$ & $\dfrac{F(\omega)}{\jj\omega} + \pi F(0)\delta(\omega)$ \\[6pt]
Convolution & $f(t) \conv g(t)$ & $F(\omega) \cdot G(\omega)$ \\[6pt]
Modulation & $f(t) \cdot g(t)$ & $\dfrac{1}{2\pi}F(\omega) \conv G(\omega)$ \\[6pt]
Parseval's Theorem & $\displaystyle\int_{-\infty}^{\infty}|f(t)|^2\,\dd t$ & $\dfrac{1}{2\pi}\displaystyle\int_{-\infty}^{\infty}|F(\omega)|^2\,\dd\omega$ \\[6pt]
Duality & $F(t)$ & $2\pi f(-\omega)$ \\
\bottomrule
\end{tabular}
\end{table}

\begin{warningbox}[title=Convolution vs. Modulation]
Note the duality between convolution and modulation:
\begin{itemize}
    \item Convolution in time $\Leftrightarrow$ Multiplication in frequency
    \item Multiplication in time $\Leftrightarrow$ Convolution in frequency (scaled by $1/2\pi$)
\end{itemize}
This duality is fundamental to understanding filtering, modulation, and signal processing.
\end{warningbox}

\subsubsection{Relationship to Laplace Transform}
For causal signals, $F(\omega) = F(s)|_{s=\jj\omega}$ when the Laplace transform converges on the imaginary axis.

\begin{keyconceptbox}[title=When Does the Fourier Transform Exist?]
The Fourier transform of $f(t)$ exists if:
\begin{enumerate}
    \item $f(t)$ is absolutely integrable: $\displaystyle\int_{-\infty}^{\infty}|f(t)|\,\dd t < \infty$, or
    \item $f(t)$ is square integrable: $\displaystyle\int_{-\infty}^{\infty}|f(t)|^2\,\dd t < \infty$
\end{enumerate}
Functions like $\sin(\omega_0 t)$ or the unit step $u(t)$ do not satisfy these conditions, but their Fourier transforms can be defined using generalized functions (distributions) involving impulses $\delta(\omega)$.
\end{keyconceptbox}

\section{Frequency Response}

For a stable LTI system with transfer function $G(s)$, the \textbf{frequency response} is:
\begin{equation}
G(\jj\omega) = |G(\jj\omega)| \ee^{\jj\angle G(\jj\omega)}
\end{equation}

\begin{itemize}
    \item $|G(\jj\omega)|$ = magnitude (gain) at frequency $\omega$
    \item $\angle G(\jj\omega)$ = phase shift at frequency $\omega$
\end{itemize}

For a sinusoidal input $u(t) = A\sin(\omega t)$, the steady-state output is:
\begin{equation}
y_{ss}(t) = A|G(\jj\omega)|\sin(\omega t + \angle G(\jj\omega))
\end{equation}

\section{Bode Plots}

Bode plots display magnitude (in dB) and phase versus frequency (on a log scale):
\begin{equation}
|G(\jj\omega)|_{\text{dB}} = 20\log_{10}|G(\jj\omega)|
\end{equation}

\begin{figure}[htbp]
\centering
\textbf{Magnitude Plot}

\begin{tikzpicture}
\begin{semilogxaxis}[
    width=12cm, height=6cm,
    xlabel={Frequency $\omega$ (rad/s)},
    ylabel={Magnitude (dB)},
    xmin=0.1, xmax=100,
    ymin=-45, ymax=25,
    grid=major,
    grid style={gray!30},
    title={\large\bfseries First-Order System: $G(s) = \frac{10}{s+1}$},
    title style={at={(0.5,-0.12)}, anchor=north},
    legend pos=south west
]

% Actual magnitude response
\addplot[UofSCAtlantic, very thick, domain=0.1:100, samples=150] {20*log10(10/sqrt(1+x^2))};
\addlegendentry{Actual Response}

% Asymptotic approximations
\addplot[UofSCGarnet, dashed, very thick, domain=0.1:0.99, samples=50] {20*log10(10)};
\addplot[UofSCGarnet, dashed, very thick, domain=1:100, samples=50] {20*log10(10) - 20*log10(x)};
\addlegendentry{Asymptotic Approximation}

% Corner frequency marker
\draw[UofSCHorseshoe, thick, dashed] (1, -45) -- (1, 25);
\node[UofSCHorseshoe, font=\small, anchor=south west] at (1, 24.5) {Corner: $\omega_c = 1$ rad/s};

% -3dB point marker
\draw[UofSCRose, thin, dotted] (0.1, 17) -- (100, 17);
\node[UofSCRose, font=\small, anchor=west] at (100.3, 17) {$-3$ dB point};

% Point on curve at corner frequency
\node[circle, fill, UofSCAtlantic, minimum size=4pt] at (1, 17) {};

\end{semilogxaxis}
\end{tikzpicture}

\vspace{0.8cm}

\textbf{Phase Plot}

\begin{tikzpicture}
\begin{semilogxaxis}[
    width=12cm, height=6cm,
    xlabel={Frequency $\omega$ (rad/s)},
    ylabel={Phase (degrees)},
    xmin=0.1, xmax=100,
    ymin=-100, ymax=10,
    grid=major,
    grid style={gray!30},
    ytick={0,-45,-90},
    yticklabels={$0^{\circ}$, $-45^{\circ}$, $-90^{\circ}$},
    legend pos=north east
]

% Phase response
\addplot[UofSCCongaree, very thick, domain=0.1:100, samples=150] {-atan(x)*180/pi};
\addlegendentry{Phase Response}

% Corner frequency marker
\draw[UofSCHorseshoe, thick, dashed] (1, -100) -- (1, 10);

% -45^{\circ} point marker
\draw[UofSCRose, thin, dotted] (0.1, -45) -- (100, -45);
\node[UofSCRose, font=\small, anchor=west] at (100.3, -45) {$-45^{\circ}$ at corner};

% Point on curve at corner frequency
\node[circle, fill, UofSCCongaree, minimum size=4pt] at (1, -45) {};

\end{semilogxaxis}
\end{tikzpicture}

\caption{Bode plot for a first-order system $G(s) = \frac{10}{s+1}$. The magnitude plot (top, Atlantic blue) shows the steady-state gain of 20 dB below the corner frequency $\omega_c = 1$ rad/s, with a -20 dB/decade slope afterward. The phase plot (bottom, Congaree teal) shows $-45^{\circ}$ phase lag at the corner frequency. Dashed Garnet lines show asymptotic approximations. The UofSCRose accent color highlights key features: the -3 dB magnitude point and the $-45^{\circ}$ phase point at the corner frequency.}
\label{fig:bode_first_order}
\end{figure}


\begin{examplebox}[title=Example: Drawing a Bode Plot from Frequency Response]
Construct the Bode plot for the first-order low-pass filter:
\[
G(s) = \frac{K}{1 + s/\omega_c}
\]
where $K = 10$ (DC gain) and $\omega_c = 100$ rad/s (corner frequency).

\textbf{Step 1: Find the frequency response}

Substitute $s = \jj\omega$:
\[
G(\jj\omega) = \frac{K}{1 + \jj\omega/\omega_c} = \frac{10}{1 + \jj\omega/100}
\]

\textbf{Step 2: Calculate magnitude}
\[
|G(\jj\omega)| = \frac{K}{\sqrt{1 + (\omega/\omega_c)^2}} = \frac{10}{\sqrt{1 + (\omega/100)^2}}
\]

In decibels:
\[
|G(\jj\omega)|_{\text{dB}} = 20\log_{10}(K) - 20\log_{10}\sqrt{1 + (\omega/\omega_c)^2} = 20 - 10\log_{10}(1 + \omega^2/10000)
\]

\textbf{Step 3: Calculate phase}
\[
\angle G(\jj\omega) = -\arctan\left(\frac{\omega}{\omega_c}\right) = -\arctan\left(\frac{\omega}{100}\right)
\]

\textbf{Step 4: Asymptotic approximations}

\underline{Magnitude:}
\begin{itemize}
    \item For $\omega \ll \omega_c$: $|G|_{\text{dB}} \approx 20\log_{10}(K) = 20$ dB (flat)
    \item For $\omega \gg \omega_c$: $|G|_{\text{dB}} \approx 20\log_{10}(K) - 20\log_{10}(\omega/\omega_c)$ (slope: $-20$ dB/decade)
    \item At $\omega = \omega_c$: $|G|_{\text{dB}} = 20 - 3 = 17$ dB (corner point, 3 dB below)
\end{itemize}

\underline{Phase:}
\begin{itemize}
    \item For $\omega \ll \omega_c/10$: $\angle G \approx 0^\circ$
    \item At $\omega = \omega_c$: $\angle G = -45^\circ$
    \item For $\omega \gg 10\omega_c$: $\angle G \approx -90^\circ$
\end{itemize}

\textbf{Result:} The Bode plot shows a flat response at low frequencies with 20 dB gain, rolling off at $-20$ dB/decade above the corner frequency, with phase shifting from $0^\circ$ to $-90^\circ$.
\end{examplebox}

\begin{keyconceptbox}[title=Standard Bode Plot Building Blocks]
Any transfer function can be decomposed into these basic elements:
\begin{enumerate}
    \item \textbf{Constant gain $K$:} Flat magnitude at $20\log_{10}|K|$ dB, phase $0^\circ$ or $\pm 180^\circ$
    \item \textbf{Poles/zeros at origin $(s)^{\pm n}$:} $\pm 20n$ dB/decade slope through 0 dB at $\omega = 1$, phase $\pm 90n^\circ$
    \item \textbf{Real pole/zero $(1 + s/\omega_c)^{\pm 1}$:} $\pm 20$ dB/decade above $\omega_c$, phase change of $\pm 90^\circ$
    \item \textbf{Complex poles/zeros:} Resonant peak/notch near $\omega_n$, phase change of $\pm 180^\circ$
\end{enumerate}
\end{keyconceptbox}

\section{Discrete Fourier Transform and FFT}

When working with sampled data, we use the Discrete Fourier Transform (DFT) instead of the continuous Fourier transform.

\begin{definitionbox}[title=Definition 5.2: Discrete Fourier Transform (DFT)]
For a sequence of $N$ samples $\{x[0], x[1], \ldots, x[N-1]\}$, the DFT is defined as:
\begin{equation}
X[k] = \sum_{n=0}^{N-1} x[n] \ee^{-\jj 2\pi kn/N}, \quad k = 0, 1, \ldots, N-1
\end{equation}

The inverse DFT (IDFT) recovers the original sequence:
\begin{equation}
x[n] = \frac{1}{N} \sum_{k=0}^{N-1} X[k] \ee^{\jj 2\pi kn/N}, \quad n = 0, 1, \ldots, N-1
\end{equation}
\end{definitionbox}

\subsection{Connection to Continuous Fourier Transform}

The DFT approximates the continuous Fourier transform at discrete frequencies. If the continuous signal $f(t)$ is sampled at rate $f_s = 1/T_s$ to produce $N$ samples:
\begin{itemize}
    \item The DFT frequency bins are spaced at $\Delta f = f_s/N$ Hz
    \item The $k$-th DFT coefficient corresponds to frequency $f_k = k \cdot f_s/N$
    \item The relationship is: $X[k] \approx T_s \cdot F(2\pi f_k)$ for band-limited signals
\end{itemize}

\textbf{Sampling and the Nyquist Frequency}

\vspace{0.5cm}

% Continuous signal with adequate sampling
\begin{tikzpicture}
\begin{axis}[
    width=13cm, height=4cm,
    xlabel={Time $t$ (seconds)},
    ylabel={Amplitude},
    xmin=0, xmax=2,
    ymin=-1.3, ymax=1.3,
    grid=major,
    grid style={gray!30},
    title={\normalsize Adequate Sampling: $f_s = 10$ Hz $> 2f_{\text{signal}} = 6$ Hz (Nyquist criterion satisfied)},
    title style={at={(0.5,1.05)}, anchor=south}
]

% Original continuous signal: 3 Hz sine wave
\addplot[UofSC50Black, dashed, thick, domain=0:2, samples=500]
    {sin(deg(2*pi*3*x))};

% Sample points at 10 Hz (every 0.1 s)
\addplot+[UofSCAtlantic, only marks, mark=*, mark size=3pt]
    coordinates {
        (0, 0)
        (0.1, 0.951)
        (0.2, 0.588)
        (0.3, -0.588)
        (0.4, -0.951)
        (0.5, 0)
        (0.6, 0.951)
        (0.7, 0.588)
        (0.8, -0.588)
        (0.9, -0.951)
        (1.0, 0)
        (1.1, 0.951)
        (1.2, 0.588)
        (1.3, -0.588)
        (1.4, -0.951)
        (1.5, 0)
        (1.6, 0.951)
        (1.7, 0.588)
        (1.8, -0.588)
        (1.9, -0.951)
        (2.0, 0)
    };

% Reconstructed signal (matches original well)
\addplot[UofSCAtlantic, very thick, domain=0:2, samples=500]
    {sin(deg(2*pi*3*x))};

% Legend entries
\addlegendentry{Original: $f = 3$ Hz}
\addlegendentry{Sample points}
\addlegendentry{Reconstructed}

% Annotation
\node[UofSCHorseshoe, font=\footnotesize, fill=white, inner sep=2pt] at (1.5, 0.85) {$f_s > 2f_{\text{max}}$ : No aliasing};

\end{axis}
\end{tikzpicture}

\vspace{0.6cm}

% Undersampled signal showing aliasing
\begin{tikzpicture}
\begin{axis}[
    width=13cm, height=4cm,
    xlabel={Time $t$ (seconds)},
    ylabel={Amplitude},
    xmin=0, xmax=2,
    ymin=-1.3, ymax=1.3,
    grid=major,
    grid style={gray!30},
    title={\normalsize Undersampling: $f_s = 4$ Hz $< 2f_{\text{signal}} = 6$ Hz (Aliasing occurs!)},
    title style={at={(0.5,1.05)}, anchor=south}
]

% Original continuous signal: 3 Hz sine wave
\addplot[UofSC50Black, dashed, thick, domain=0:2, samples=500]
    {sin(deg(2*pi*3*x))};

% Sample points at 4 Hz (every 0.25 s) - sampling a 3 Hz signal
% sin(2*pi*3*t) at t = 0, 0.25, 0.5, 0.75, 1.0, ...
\addplot+[UofSCGarnet, only marks, mark=*, mark size=3pt]
    coordinates {
        (0, 0)
        (0.25, -0.707)
        (0.5, 0)
        (0.75, 0.707)
        (1.0, 0)
        (1.25, -0.707)
        (1.5, 0)
        (1.75, 0.707)
        (2.0, 0)
    };

% Aliased signal: appears as 1 Hz signal (|3 - 4| = 1 Hz alias)
\addplot[UofSCRose, very thick, domain=0:2, samples=300]
    {-sin(deg(2*pi*1*x))};

% Legend
\addlegendentry{Original: $f = 3$ Hz}
\addlegendentry{Sample points}
\addlegendentry{Aliased: $f_{\text{alias}} = |f - f_s| = 1$ Hz}

% Annotation
\node[UofSCGarnet, font=\footnotesize, fill=white, inner sep=2pt] at (1.5, -0.85) {$f_s < 2f_{\text{max}}$ : \textbf{Aliasing!}};

\end{axis}
\end{tikzpicture}

\vspace{0.6cm}

% Frequency domain illustration
\begin{tikzpicture}
\begin{axis}[
    width=6.2cm, height=4.5cm,
    xlabel={Frequency $f$ (Hz)},
    ylabel={$|X(f)|$},
    xmin=-8, xmax=8,
    ymin=0, ymax=1.3,
    grid=major,
    grid style={gray!30},
    title={\normalsize Adequate Sampling Spectrum},
    title style={at={(0.5,1.05)}, anchor=south},
    xtick={-6,-4,-3,-2,0,2,3,4,6},
    ytick={0, 0.5, 1}
]

% Original spectrum (impulses at +/- 3 Hz)
\addplot+[UofSCAtlantic, very thick, ycomb, mark=triangle*, mark options={fill=UofSCAtlantic, scale=1.5}]
    coordinates {(-3, 1) (3, 1)};

% Nyquist frequency markers
\draw[UofSCHorseshoe, dashed, very thick] (-5, 0) -- (-5, 1.2);
\draw[UofSCHorseshoe, dashed, very thick] (5, 0) -- (5, 1.2);

% Labels
\node[UofSCAtlantic, font=\footnotesize, anchor=south] at (3, 1.02) {$f_0$};
\node[UofSCAtlantic, font=\footnotesize, anchor=south] at (-3, 1.02) {$-f_0$};
\node[UofSCHorseshoe, font=\footnotesize, anchor=south] at (5, 1.22) {$f_s/2$};
\node[UofSCHorseshoe, font=\footnotesize, anchor=south] at (-5, 1.22) {$-f_s/2$};

% Shade the Nyquist band
\fill[UofSCHorseshoe, opacity=0.1] (-5, 0) rectangle (5, 1.2);

\end{axis}
\end{tikzpicture}
\hfill
\begin{tikzpicture}
\begin{axis}[
    width=6.2cm, height=4.5cm,
    xlabel={Frequency $f$ (Hz)},
    ylabel={$|X(f)|$},
    xmin=-8, xmax=8,
    ymin=0, ymax=1.3,
    grid=major,
    grid style={gray!30},
    title={\normalsize Undersampled Spectrum (Aliasing)},
    title style={at={(0.5,1.05)}, anchor=south},
    xtick={-6,-4,-3,-2,-1,0,1,2,3,4,6},
    ytick={0, 0.5, 1}
]

% Original spectrum (impulses at +/- 3 Hz) - shown faded
\addplot+[UofSC50Black, thick, ycomb, mark=triangle*, mark options={fill=UofSC50Black, scale=1.2}]
    coordinates {(-3, 0.7) (3, 0.7)};

% Aliased frequencies (folded into Nyquist band)
% With fs=4, f_N=2, signal at 3 Hz aliases to |3-4|=1 Hz
\addplot+[UofSCRose, very thick, ycomb, mark=*, mark options={fill=UofSCRose, scale=1.5}]
    coordinates {(-1, 1) (1, 1)};

% Nyquist frequency markers
\draw[UofSCGarnet, dashed, very thick] (-2, 0) -- (-2, 1.2);
\draw[UofSCGarnet, dashed, very thick] (2, 0) -- (2, 1.2);

% Labels
\node[UofSC50Black, font=\footnotesize, anchor=south] at (3, 0.72) {$f_0$};
\node[UofSCRose, font=\footnotesize, anchor=south] at (1, 1.02) {$f_{\text{alias}}$};
\node[UofSCGarnet, font=\footnotesize, anchor=south] at (2, 1.22) {$f_s/2$};
\node[UofSCGarnet, font=\footnotesize, anchor=south] at (-2, 1.22) {$-f_s/2$};

% Shade the Nyquist band
\fill[UofSCGarnet, opacity=0.1] (-2, 0) rectangle (2, 1.2);

% Arrow showing folding
\draw[UofSCRose, thick, ->, >=stealth] (3, 0.5) to[out=180, in=0] (1, 0.7);
\node[UofSCRose, font=\tiny] at (2.3, 0.3) {Folds};

\end{axis}
\end{tikzpicture}

\vspace{0.4cm}

% Nyquist theorem statement box
\begin{tikzpicture}
\node[draw=UofSCGarnet, thick, fill=UofSCGarnet!8, minimum width=12.5cm, minimum height=1.8cm, sharp corners, inner sep=8pt, align=center] {
    \textbf{Nyquist-Shannon Sampling Theorem:} A bandlimited signal with maximum frequency $f_{\text{max}}$ \\[2pt]
    can be perfectly reconstructed from samples if $f_s > 2f_{\text{max}}$. \\[4pt]
    \textcolor{UofSCGarnet}{\textbf{Nyquist Frequency:}} $f_N = f_s/2$ \quad \textcolor{UofSCGarnet}{\textbf{Nyquist Rate:}} $f_s = 2f_{\text{max}}$
};
\end{tikzpicture}



\begin{warningbox}[title=Nyquist Sampling Theorem]
To avoid aliasing, the sampling frequency must satisfy:
\[
f_s > 2 f_{\max}
\]
where $f_{\max}$ is the highest frequency component in the signal. The frequency $f_s/2$ is called the \textbf{Nyquist frequency}. Frequency components above the Nyquist frequency will ``fold back'' and appear as lower frequencies (aliasing).
\end{warningbox}

\subsection{Fast Fourier Transform (FFT)}

Direct computation of the DFT requires $\mathcal{O}(N^2)$ operations. The \textbf{Fast Fourier Transform (FFT)} is an efficient algorithm that computes the DFT in $\mathcal{O}(N \log N)$ operations by exploiting symmetry and periodicity of the complex exponentials.

\begin{keyconceptbox}[title=FFT Key Points]
\begin{itemize}
    \item The FFT computes the \textit{exact same result} as the DFT---it is not an approximation
    \item Most efficient when $N$ is a power of 2 (e.g., 256, 512, 1024, 2048)
    \item Speedup factor: For $N = 1024$, FFT is about 100$\times$ faster than direct DFT
    \item The Cooley-Tukey algorithm (1965) is the most common FFT implementation
    \item Modern libraries (FFTW, NumPy, MATLAB) use highly optimized FFT routines
\end{itemize}
\end{keyconceptbox}

\begin{practicalbox}[title=Using FFT in Practice]
When using FFT for spectral analysis:
\begin{enumerate}
    \item \textbf{Zero-padding:} Add zeros to increase frequency resolution (interpolation in frequency domain)
    \item \textbf{Windowing:} Apply a window function (Hanning, Hamming, Blackman) to reduce spectral leakage
    \item \textbf{Averaging:} Use Welch's method (overlapping windowed segments) for noise reduction
    \item \textbf{Scaling:} Normalize properly to get correct magnitude units
\end{enumerate}

In MATLAB/Python:
\begin{lstlisting}[language=Python]
import numpy as np
from scipy.fft import fft, fftfreq

# Sample signal
fs = 1000  # Sampling frequency (Hz)
t = np.arange(0, 1, 1/fs)
x = np.sin(2*np.pi*50*t) + 0.5*np.sin(2*np.pi*120*t)

# Compute FFT
N = len(x)
X = fft(x)
freq = fftfreq(N, 1/fs)

# Magnitude spectrum (single-sided)
X_mag = 2*np.abs(X[:N//2])/N
\end{lstlisting}
\end{practicalbox}

% ============================================================================
